% Bab 7: Beta Testing (User Acceptance Testing)
\chapter{Beta Testing (User Acceptance Testing)}
\label{bab:beta-testing}

\section{Responden}

\begin{longtable}{|L{0.30}|L{0.70}|}
\caption{Profil Responden Beta Testing}\\
\hline
\cellcolor[HTML]{D9EAF7}\textbf{Atribut} & \cellcolor[HTML]{D9EAF7}\textbf{Keterangan} \\ \hline
\endfirsthead
\hline
\cellcolor[HTML]{D9EAF7}\textbf{Atribut} & \cellcolor[HTML]{D9EAF7}\textbf{Keterangan} \\ \hline
\endhead
Target pengguna & Pengunjung kawasan kuliner Sudirman Street Food (lokal/non-lokal) dan pengunjung dengan waktu terbatas \\ \hline
Jumlah responden & [isi jumlah responden; minimal 2 orang di luar tim pengembang] \\ \hline
Cara rekrutmen & [isi cara rekrutmen] \\ \hline
Lokasi dan waktu & Sudirman Street Day and Night Market, [isi tanggal \& waktu] \\ \hline
\end{longtable}

\textit{Jumlah sampel minimal 2 orang di luar tim pengembang.}

\section{Tugas Penggunaan}

\begin{longtable}{|C{0.08}|L{0.50}|C{0.20}|L{0.22}|}
\caption{Daftar Tugas Penggunaan Beta Testing}\\
\hline
\cellcolor[HTML]{D9EAF7}\textbf{Tugas} & \cellcolor[HTML]{D9EAF7}\textbf{Instruksi untuk Responden} & \cellcolor[HTML]{D9EAF7}\textbf{Berhasil tanpa bantuan} & \cellcolor[HTML]{D9EAF7}\textbf{Catatan} \\ \hline
\endfirsthead
\hline
\cellcolor[HTML]{D9EAF7}\textbf{Tugas} & \cellcolor[HTML]{D9EAF7}\textbf{Instruksi untuk Responden} & \cellcolor[HTML]{D9EAF7}\textbf{Berhasil tanpa bantuan} & \cellcolor[HTML]{D9EAF7}\textbf{Catatan} \\ \hline
\endhead
T1 & Buka aplikasi untuk pertama kali, masuk/daftar akun lokal, ikuti onboarding (izinkan lokasi dan notifikasi, pilih preferensi rasa serta opsi halal bila perlu), lalu aktifkan Mode Jelajah. & \_\_ dari \_\_ & \\ \hline
T2 & Setelah memasuki kawasan, rencanakan trail untuk 60 menit dengan budget Rp50.000, lalu mulai trail. & \_\_ dari \_\_ & \\ \hline
T3 & Ikuti navigasi menuju tenant pertama dengan HP di saku, lalu dengarkan narasi tenant tersebut. & \_\_ dari \_\_ & \\ \hline
T4 & Cari tenant melalui kolom pencarian, terapkan filter harga <Rp15rb, lalu jadikan salah satu tenant sebagai tujuan. & \_\_ dari \_\_ & \\ \hline
T5 & Setelah selesai makan di tenant, lihat rekomendasi tenant berikutnya, lalu akhiri sesi dengan Selesai Jelajah. & \_\_ dari \_\_ & \\ \hline
\end{longtable}

\section{Instrumen Kuesioner}

\textit{Gunakan skala Likert: 1 = Sangat Tidak Setuju (STS), 2 = Tidak Setuju (TS), 3 = Netral (N), 4 = Setuju (S), 5 = Sangat Setuju (SS).}

\begin{longtable}{|L{0.22}|C{0.10}|L{0.68}|}
\caption{Instrumen Kuesioner Skala Likert}\\
\hline
\cellcolor[HTML]{D9EAF7}\textbf{Aspek} & \cellcolor[HTML]{D9EAF7}\textbf{Kode} & \cellcolor[HTML]{D9EAF7}\textbf{Pertanyaan} \\ \hline
\endfirsthead
\hline
\cellcolor[HTML]{D9EAF7}\textbf{Aspek} & \cellcolor[HTML]{D9EAF7}\textbf{Kode} & \cellcolor[HTML]{D9EAF7}\textbf{Pertanyaan} \\ \hline
\endhead
Kualitas Sistem & KS1 & Aplikasi mudah dipahami saat pertama kali digunakan. \\ \hline
 & KS2 & Notifikasi POI muncul pada saat yang tepat ketika saya mendekati tenant. \\ \hline
 & KS3 & Mode Jelajah dan pemicu lokasi berjalan stabil tanpa pemicu yang muncul berulang. \\ \hline
Kualitas Informasi & KI1 & Informasi tenant (menu, harga, dan cerita budaya) yang ditampilkan jelas dan mudah dipahami. \\ \hline
 & KI2 & Narasi audio terdengar jelas dan mudah diikuti. \\ \hline
 & KI3 & Rekomendasi tenant berikutnya disertai alasan yang masuk akal. \\ \hline
Kualitas Layanan & KL1 & Navigasi (rute, panah, jarak, dan panduan audio) membantu saya menemukan tenant yang ingin dikunjungi. \\ \hline
 & KL2 & Trail yang disusun aplikasi sesuai dengan waktu dan budget yang saya tentukan. \\ \hline
 & KL3 & Pencarian dan filter membantu saya menemukan tenant yang sesuai. \\ \hline
Kepuasan Pengguna & KP1 & Secara keseluruhan saya puas menggunakan aplikasi ini. \\ \hline
 & KP2 & Saya akan menggunakan aplikasi ini lagi saat berkunjung ke kawasan kuliner. \\ \hline
 & KP3 & Aplikasi ini membuat eksplorasi kuliner saya lebih hemat waktu. \\ \hline
\end{longtable}

\section{Hasil Kuesioner per Pertanyaan}

\textit{Isi jumlah responden untuk setiap pilihan jawaban. Total Skor = (1 \(\times\) jumlah STS) + (2 \(\times\) jumlah TS) + (3 \(\times\) jumlah N) + (4 \(\times\) jumlah S) + (5 \(\times\) jumlah SS).}

\begin{longtable}{|C{0.09}|L{0.45}|C{0.06}|C{0.06}|C{0.06}|C{0.06}|C{0.06}|C{0.16}|}
\caption{Hasil Kuesioner Aspek}\\
\hline
\cellcolor[HTML]{D9EAF7}\textbf{Kode} & \cellcolor[HTML]{D9EAF7}\textbf{Pertanyaan} & \cellcolor[HTML]{D9EAF7}\textbf{STS} & \cellcolor[HTML]{D9EAF7}\textbf{TS} & \cellcolor[HTML]{D9EAF7}\textbf{N} & \cellcolor[HTML]{D9EAF7}\textbf{S} & \cellcolor[HTML]{D9EAF7}\textbf{SS} & \cellcolor[HTML]{D9EAF7}\textbf{Total Skor} \\ \hline
\endfirsthead
\hline
\cellcolor[HTML]{D9EAF7}\textbf{Kode} & \cellcolor[HTML]{D9EAF7}\textbf{Pertanyaan} & \cellcolor[HTML]{D9EAF7}\textbf{STS} & \cellcolor[HTML]{D9EAF7}\textbf{TS} & \cellcolor[HTML]{D9EAF7}\textbf{N} & \cellcolor[HTML]{D9EAF7}\textbf{S} & \cellcolor[HTML]{D9EAF7}\textbf{SS} & \cellcolor[HTML]{D9EAF7}\textbf{Total Skor} \\ \hline
\endhead
KS1 & Aplikasi mudah dipahami saat pertama kali digunakan. & & & & & & \\ \hline
KS2 & Notifikasi POI muncul pada saat yang tepat ketika saya mendekati tenant. & & & & & & \\ \hline
KS3 & Mode Jelajah dan pemicu lokasi berjalan stabil tanpa pemicu yang muncul berulang. & & & & & & \\ \hline
KI1 & Informasi tenant (menu, harga, dan cerita budaya) yang ditampilkan jelas dan mudah dipahami. & & & & & & \\ \hline
KI2 & Narasi audio terdengar jelas dan mudah diikuti. & & & & & & \\ \hline
KI3 & Rekomendasi tenant berikutnya disertai alasan yang masuk akal. & & & & & & \\ \hline
KL1 & Navigasi (rute, panah, jarak, dan panduan audio) membantu saya menemukan tenant yang ingin dikunjungi. & & & & & & \\ \hline
KL2 & Trail yang disusun aplikasi sesuai dengan waktu dan budget yang saya tentukan. & & & & & & \\ \hline
KL3 & Pencarian dan filter membantu saya menemukan tenant yang sesuai. & & & & & & \\ \hline
KP1 & Secara keseluruhan saya puas menggunakan aplikasi ini. & & & & & & \\ \hline
KP2 & Saya akan menggunakan aplikasi ini lagi saat berkunjung ke kawasan kuliner. & & & & & & \\ \hline
KP3 & Aplikasi ini membuat eksplorasi kuliner saya lebih hemat waktu. & & & & & & \\ \hline
\end{longtable}

\section{Perhitungan Indeks Penerimaan}

\begin{itemize}
  \item Nilai Maksimal = skala terbesar \(\times\) jumlah responden (contoh: \(5 \times 20 = 100\)).
  \item Rerata Skor aspek = \((\mbox{Total Skor pertanyaan } 1 + \dots + \mbox{Total Skor pertanyaan } N) \div N\), dengan \(N\) = jumlah pertanyaan dalam aspek.
  \item Indeks Penerimaan = \((\mbox{Rerata Skor} \div \mbox{Nilai Maksimal}) \times 100\%\).
\end{itemize}

\textit{\textbf{Contoh (aspek Kualitas Sistem, 20 responden):} Rerata Skor = \((88 + 82) \div 2 = 85\); Indeks Penerimaan = \((85 \div 100) \times 100\% = 85\%\), kategori Sangat Diterima.}

\begin{longtable}{|C{0.40}|L{0.60}|}
\caption{Kategori Indeks Penerimaan}\\
\hline
\cellcolor[HTML]{D9EAF7}\textbf{Indeks Penerimaan} & \cellcolor[HTML]{D9EAF7}\textbf{Kategori} \\ \hline
\endfirsthead
\hline
\cellcolor[HTML]{D9EAF7}\textbf{Indeks Penerimaan} & \cellcolor[HTML]{D9EAF7}\textbf{Kategori} \\ \hline
\endhead
0\% -- 19,99\% & Sangat Tidak Diterima \\ \hline
20\% -- 39,99\% & Tidak Diterima \\ \hline
40\% -- 59,99\% & Netral \\ \hline
60\% -- 79,99\% & Diterima \\ \hline
80\% -- 100\% & Sangat Diterima \\ \hline
\end{longtable}

\begin{longtable}{|L{0.26}|C{0.18}|C{0.18}|C{0.15}|L{0.23}|}
\caption{Rekapitulasi Indeks Penerimaan}\\
\hline
\cellcolor[HTML]{D9EAF7}\textbf{Aspek} & \cellcolor[HTML]{D9EAF7}\textbf{Rerata Skor} & \cellcolor[HTML]{D9EAF7}\textbf{Nilai Maksimal} & \cellcolor[HTML]{D9EAF7}\textbf{Indeks} & \cellcolor[HTML]{D9EAF7}\textbf{Kategori} \\ \hline
\endfirsthead
\hline
\cellcolor[HTML]{D9EAF7}\textbf{Aspek} & \cellcolor[HTML]{D9EAF7}\textbf{Rerata Skor} & \cellcolor[HTML]{D9EAF7}\textbf{Nilai Maksimal} & \cellcolor[HTML]{D9EAF7}\textbf{Indeks} & \cellcolor[HTML]{D9EAF7}\textbf{Kategori} \\ \hline
\endhead
Kualitas Sistem & & & & \\ \hline
Kualitas Informasi & & & & \\ \hline
Kualitas Layanan & & & & \\ \hline
Kepuasan Pengguna & & & & \\ \hline
\textbf{Keseluruhan} & & & & \\ \hline
\end{longtable}

\textit{Indeks keseluruhan dapat dihitung dari rata-rata Rerata Skor seluruh pertanyaan dibagi Nilai Maksimal.}

\section{Masukan Terbuka dari Responden}

\begin{longtable}{|L{0.20}|L{0.42}|L{0.38}|}
\caption{Masukan Terbuka dari Responden}\\
\hline
\cellcolor[HTML]{D9EAF7}\textbf{Responden} & \cellcolor[HTML]{D9EAF7}\textbf{Masukan} & \cellcolor[HTML]{D9EAF7}\textbf{Tindak Lanjut} \\ \hline
\endfirsthead
\hline
\cellcolor[HTML]{D9EAF7}\textbf{Responden} & \cellcolor[HTML]{D9EAF7}\textbf{Masukan} & \cellcolor[HTML]{D9EAF7}\textbf{Tindak Lanjut} \\ \hline
\endhead
 & & \\ \hline
 & & \\ \hline
 & & \\ \hline
 & & \\ \hline
\end{longtable}
